\subsection{An Express Exchange}

An express send needs both ends configured: the sender uses an express
evaluator, and the receiver uses an express metaobject generator.
Begin with the original clock while it is idle. In a fresh session at
the repository root, load its library and definitions:

\begin{codeblock}
\begin{lstlisting}
(inc-lib "lib/express" :use-source t)
(inc "examples/misc/express-ex" :use-source t)
\end{lstlisting}
\end{codeblock}

The example defines the group \code{express} and its clock.
Place the client in that group too, so its evaluator understands the
extra \code{:express} send argument:

\begin{codeblock}
\begin{lstlisting}
[object finite-express-sender
 (group express)
 (script
  (=> :peek ![clock <== :peek :express])
  (=> :past-peek @ R [clock <= :peek @ R :express])
  (=> :atomic-result !(atomic 42)))]
\end{lstlisting}
\end{codeblock}

Try these short exchanges:

\begin{codeblock}
\begin{lstlisting}
[clock <== :peek]                         ; => 0
[finite-express-sender <== :peek]          ; => 0
[finite-express-sender <== :past-peek]     ; => 0
[finite-express-sender <== :atomic-result] ; => 42
\end{lstlisting}
\end{codeblock}

The first query is ordinary. The second goes through an express-capable
sender. The third forwards the caller's reply destination through a
past-type express send; \code{@ R}, not \code{<=} alone, supplies that
reply path. The last expression shows that \code{atomic} returns its
body's value.

This exercise studies routing to an idle receiver, not interruption of
a running computation. Section~\ref{sec:original-express} follows the
original clock's scripts, mode transitions, and evaluator front end.
